\begin{abstract}
The interaction field of the SHOW3D benchmark assigns to each of the 21 joints of a hand the three-dimensional vector to the nearest point of the manipulated object, observed from the two cameras of an egocentric headset. Published solutions regress these 126 numbers directly, although the field is a deterministic function of three supervised quantities: the hand joints, the object pose and a known object mesh. We propose factored interaction fields (\method{}), a transformer decoder over frozen DINOv3 tokens that predicts camera-frame joints and the object pose, evaluates the field through a differentiable soft nearest-vertex operator on the posed mesh, and fuses this geometric estimate with a direct regression head by predicted precision. On held-out subjects, \method{} reduces the error of a matched direct regressor from 46.2 to 44.4\,mm and from 35.7 to 34.1\,mm on a second subject pair, against 60.5\,mm for the official baseline; a single-view \method{} outperforms the two-view direct model. The gain grows with difficulty, from 0.6\,mm for joints within 15\,mm of the object to 21\,mm when the hand leaves the field of view. Inside the model the improvement is carried by the direct head, whose representation is shaped by the geometric supervision; inference-time fusion contributes 0.1\,mm and the geometric head's predicted scale saturated in four of five runs. The advantage does not carry over to HOT3D in four cross-dataset settings, which we trace to the object-pose branch. We release code, derived labels and the HOT3D-IF cross-dataset protocol.
\end{abstract}
